\scanpage{21}
% archive foliation: 20
\begin{center}- 89 -\end{center}\hfill n\textsuperscript{o} 268

et si $f$, $f'$ sont deux fonctions telles que $(f)=X-X_0$;

et $(f')=X'-X_0$, il faut montrer que pour toute constante $c$, le diviseur $X_c$ tel que $(f+cf')=X_c-X_0$ est dans $\mathcal L_N$. Soit $g$ une fonction et $Y$ un diviseur tels que $X_0=(g)+Y$, et que $P\notin\operatorname{supp}(Y)$. Alors on a $(fg)=X-Y$ et $(f'g)=X'-Y$.

Ceci montre que $fg$ et $f'g$ sont dans $I_N$, donc que $(f+cf')g$ est dans $I_N$ On a évidemment
\[
(f+cf')g=X_c-Y.
\]
et par définition, on voit que $X_c$ satisfait aux conditions $I_N$. Ceci démontre que $\mathcal L_N$ est un système linéaire. Pour $N$ suffisamment grand, on voit tout de suite que $\mathcal L_N$ est vide. Quant au fait que la dimension baisse au plus d'une unité quand $N$ augmente.d'une unité, c'est clair.

\uline{Proposition} - Soit $\mathcal L$ un système linéaire \uline{complet}, et $\{I_N\}$ un système de conditions linéaires en un point simple $P$ de $V$. Alors :

(i) Il existe un entier $N(\geq-1)$ uniquement déterminé par la propriété suivante : Au moins un diviseur $X$ de $\mathcal L$ satisfait les condition $I_N$ et aucun diviseur de $\mathcal L$ ne satisfait aux conditions $I_{N+1}$. Pour cet entier $N$, il existe un diviseur $X$ et un seul de $\mathcal L$ satisfaisant $I_N$.

(ii) Si $V$ est définie sur $k$, $P$ rationnel sur un corps $K\supset k$, et s'il existe un diviseur $D$ sur $V$, rationnel sur $K$, linéairement équivalent aux diviseurs de $\mathcal L$, alors le diviseur unique $X$ du (i) est rationnel sur $K$.

\uline{Démonstration} - Le (i) est une conséquence immédiate de la proposition qui précède. Passons au (ii). En employant le fait que tout diviseur est localement linéairement équivalent à $0$ en un point simple,
